\subsection{All odd terminals with nonunit shortest arm}\label{proof:R2-C020}
Let the original residual list be
$(\ell,s_1,t_1,\ldots,s_j,t_j)$, with
$2\le\ell<s_1<t_1<\cdots<s_j<t_j$, $s_1\ge4$, and finite $j\ge1$.
Fix the actual terminal $F=\bar f_\ell\prod_{i=1}^jh_{s_i}f_{t_i}$,
orient all connector edges toward the receiving root, and orient the
receiving arm outward. We prove, for every independent $a\ge3,L\ge1$,
\begin{equation}\tag{R20.1.1}\label{eq:R20:claim}
\E[f_a(T^LF)^4]\le1-d/5.
\end{equation}
Together with R2-C018, whose second shortest arm is three,
this covers every odd residual list with at least three entries and
shortest entry at least two. Actual equal-arm decorations can be restored in place.

Use the original finite receiving notation, exact covariance
\eqref{eq:R17:covariance}, and centered Taylor identity
\eqref{eq:R17:taylor}. In particular
$\E C_a\le1-57d/32$, $C_a\le(1-d)p+(6/5)\sqrt d$,
$\sigma(C_a)\le(1+d)\sqrt d$.
The cubic inequality \eqref{eq:R12:cubic} gives the following
continuous coefficients for $0\le m=\E Y\le1$:
\[
\begin{array}{c|ccc}
C_a+Y+m\text{ upper bound}&4&7/2&18/5\\\hline
\text{curvature coefficient}&683/54<13&248/27<19/2&2453/250<10.
\end{array}
\]
In the low range $0<d\le1/8$, the exponential payment
\eqref{eq:R14:exponential} will use
\begin{equation}\tag{R20.2.4}\label{eq:R20:payment}
\begin{array}{c|ccc|c}
 &A&B&C&\text{deficit lower bound}\\\hline
j=1&39/2&53/3&64/5&(11/54)d>d/5\\
j\ge2&29&22&128/5&(7/9)d>d/5.
\end{array}
\end{equation}
Indeed $C\ge A/2$, $Ad/2\le39/32$ or $29/16$, both below two,
and $e^2<9$.

Here is the exact finite anchored input from Q5-C181, Section 3.
Let $R$ be the sum of the disjoint original power intervals of all
actual $h_sf_t$ pairs, and let $\Gamma$ be its actual loss.
Put $\mu_\ell=\E\bar f_\ell\in[0,1]$, $D_i=\E B^iR$.
Then $F\le\bar f_\ell e^{R-\Gamma}$, $|R|^2<d/3$, and
\begin{equation}\tag{R20.3.table}\label{eq:R20:rows}
\begin{array}{c|ccccc}
\text{row}&\Gamma/d\text{ lower bound}&\rho&c_0&c_1&c_2\\\hline
I:\ j=1,s_1=4&69/20&1/5&72/125&1/6&7/100\\
II:\ j=1,s_1\ge5&16/5&3/20&3/5&1/30&1/12\\
III:\ j\ge2,s_1\ge4&16j/5&1/4&5/6&1/5&1/12
\end{array}
\end{equation}
Here $\|R\|_2<\rho\sqrt d$,
$\delta_+/d\le c_0$, $(D_1)_+/d\le c_1$, and
$(D_i)_+/d\le c_2$ for $i\ge2$. All bounds hold for every finite
mask and the whole $0<d\le1/3$.
The actual first row and finite propagation are
\begin{equation}\tag{R20.3.1}\label{eq:R20:propagation}
\begin{gathered}
T(\bar f_\ell e^R)\le\bar f_{\ell+1}+\mu_\ell BR+\delta,\\
e^\Gamma T^LF\le\bar f_n+\mu_\ell B^LR+\delta\bar f_{L-1}
+\mu_\ell\sum_{i=1}^{L-1}D_i\bar f_{L-1-i},\quad n=\ell+L.
\end{gathered}
\end{equation}
$\delta$ may be negative. In positive propagation we may replace it
by $D=c_0d\ge\delta$, and each $D_i$ only by its declared actual
positive-part bound. Resulting upper bounds $Y$ remain nonnegative.
In row I the original $\ell$ is two or three, so $\beta_{\ell+1}=d$.

For $n=3,4,5$ the correct finite bare recurrences are
\[
\begin{array}{c|c|c}
n&\bar f_n&\text{mean upper bound}\\\hline
3&q-d+B(u+Bu)&1-2d\\
4&(1-d)q-d+BG_3u&1-5d/2\\
5&(1-2d)q-\tau+B(G_4-d)u&1-11d/4.
\end{array}
\]
Their variances are at most $d$ and, uniformly,
$\bar f_n\le q+(6/5)\sqrt d-\kappa d$.
The exact short covariance formulas retained in Modules R2-C012,
R2-C014 and R2-C017 give
\[
\operatorname{Cov}(C_a,\bar f_n)\le d/12\ (n=3,4,5),\qquad
\operatorname{Cov}(C_a,\bar f_5)\le d/32.
\]
For three steps use nonpositivity at $a\ge4$ and
$d(55/1024-d)$ at $a=3$; for four steps use zero at $a=3$ and
$-57d/512+45d^2/32\le33d/512<d/12$ otherwise.
Two favorable short covariances, directly from
\eqref{eq:R17:covariance}, are
\begin{equation}\tag{R20.4.3}\label{eq:R20:shortcov}
\operatorname{Cov}(C_a,q)\le-\tau+d^2<0,\qquad
\operatorname{Cov}(C_a,\bar f_2)\le-3d/16+\theta d^2<0.
\end{equation}
For the second use $Q=1+z$, $\operatorname{Re}(G_aQ)\ge1-2x\ge3/16$,
and scalar correction at most $d^2-\tau d=d\delta_2$.
Also $\sigma(\bar f_2)^2=d+\delta_2-d^2\le(36/25)d$.

Suppose $n\le5,L=1$. Then
$Y=\bar f_n+\mu_\ell BR+D$ has
\[
\begin{gathered}
m\le1-kd,\quad k=2-c_0-c_1,\quad \operatorname{Var}Y\le(1+\sqrt\theta\rho)^2d,\\
\operatorname{Cov}(C_a,Y)\le[1/12+(9/8)\sqrt\theta\rho]d.
\end{gathered}
\]
The actual nonzero mean of $BR$ is paid by $D_1$.
The pointwise row bound gives
$Y\le q+(6/5+\rho)\sqrt d+(c_0-\kappa)d$.
All rows have $c_0-\kappa-k<0$, so
$C_a+Y+m\le3+(12/5+\rho)\sqrt d<4$; the largest increment is
$159/160<1$ because $\rho\le1/4$, $\sqrt d<3/8$.
For single pairs, $\sqrt\theta\rho<2/15$, covariance is at most
$7d/30$, and the remainder coefficient is
$13(17/15)^2+4(7/30)=3967/225<53/3$.
The mean exponent $4\Gamma/d+57/32+4k$ has $k_I=943/750$,
$k_{II}=41/30$ and exceeds $39/2$ (row II gives $9623/480$).
For multiple pairs $k=29/30$, $\sqrt\theta\rho<1/6$,
the remainder is at most $13(7/6)^2+4(13/48)=169/9<22$,
and the mean exponent is $14999/480>29$.
Thus \eqref{eq:R20:payment} applies.

If $n\le5,L=2$, then $n=4,5$, $\ell=2,3$, and use
$Y=\bar f_n+\mu_\ell B^2R+Dq+\mu_\ell(D_1)_+$.
Here $m\le1-kd$, $k=5/2-c_0-c_1-c_2$:
$k_I=2531/1500>5/3$, $k_{II}=107/60>5/3$,
$k_{III}=83/60$.
The standard deviation is at most
$[1+\theta\rho+c_0d]\sqrt d$, bounded by
$7\sqrt d/6$, $23\sqrt d/20$, and $5\sqrt d/4$ in the three rows.
For row I use $\theta/5<1/12$, $(72/125)d\le9/125$;
for II use $3\theta/20<1/16$; for III use
$\theta/4<5/48$, $(5/6)d\le5/48$.
By \eqref{eq:R20:shortcov}, covariance is at most $d/5$ for single
pairs and $d/4$ for multiple pairs.
The stronger short bound gives
$Y\le(1-d+D)q+(1+\sqrt\theta\rho)\sqrt d+(c_1-\kappa)d$.
For single pairs, $c_0\le3/5$, $c_1\le1/6$, $k\ge5/3$ give
$C_a+Y+m\le3-(463/160)d+(7/3)\sqrt d\le7/2$;
the full increment maximum is $1960/4167<1/2$.
For multiple pairs the $d$ coefficient remains negative and the
$\sqrt d$ coefficient is $71/30$, giving sum below four.
The remainder is at most
$(19/2)(49/36)+4/5=4943/360<14<53/3$ for single pairs,
and $13(25/16)+1=341/16<22$ for multiple pairs.
The mean exponents are at least
$64/5+57/32+20/3>39/2$ and
$128/5+57/32+83/15>29$, respectively.

The only remaining $n\le5,L=3$ case is $n=5,\ell=2$. Take
$Y=\bar f_5+\mu_\ell B^3R+D\bar f_2
+\mu_\ell(D_1)_+q+\mu_\ell(D_2)_+$.
Then $m\le1-kd$, $k=11/4-c_0-c_1-2c_2$:
$k_I=2801/1500>11/6$, $k_{II}=39/20>11/6$,
$k_{III}=31/20$.
Its standard deviation is at most
$[1+\theta^{3/2}\rho+(6/5)c_0d+c_1d]\sqrt d$, bounded by the
same three row values $7/6,23/20,5/4$.
The respective increments are below
$4/75+7/80+1/48<1/6$,
$1/25+9/100+1/240<3/20$, and
$1/15+1/8+1/40<1/4$.
The added $\bar f_2,q$ terms have nonpositive covariance.
The remaining tail covariance costs at most
$(9/8)\theta^{3/2}\rho d\le3d/40$; with $d/32$ use total $d/8$.
For single pairs, $D+(D_1)_+\le(23/30)d<d$ and the actual five-step
bound gives $C_a+Y+m\le3-4d+(7/3)\sqrt d\le7/2$.
Here the negative coefficient is at least
$2+\kappa+11/6-1/12>4$, the square-root coefficient is below
$11/5+\theta/5<7/3$, and the increment maximum is $49/144<1/2$.
For multiple pairs the $q$ coefficient is at most $1-(29/30)d$,
$-\tau+(D_2)_+<0$, so
$C_a+Y+m\le3+(11/5+5/48)\sqrt d<4$.
The remainders are below $14$ and at most
$13(25/16)+1/2<22$; the mean exponents are at least
$64/5+57/32+22/3>39/2$ and
$128/5+57/32+31/5>29$. This exhausts $n\le5$.

For every finite $n\ge6$, rewrite the actual propagated upper bound as
\[
Y_{\rm old}=1+G_n(B)u+\mu_\ell B^LR
-\sum_{r=3}^n\gamma_r\bar f_{n-r},
\quad
\gamma_r=\begin{cases}
\beta_r,&r<\ell+1,\\
\beta_r-D,&r=\ell+1,\\
\beta_r-\mu_\ell D_{r-\ell-1},&r>\ell+1.
\end{cases}
\]
Use the row parameters
\[
\begin{array}{c|cc}I&h=2d/5&E=0\\
II&h=d/2&E=81d/80\\III&h=3d/8&E=59d/48.\end{array}
\]
Except for the one index $r=\ell+1$, all $\gamma_r\ge h$ and
the actual mean coefficient $\beta_{n+1}-\mu_\ell D_L\ge h$.
For I, $\beta_{\ell+1}=d$ gives
$d-D\ge53d/125>2d/5$, and the other tight coefficient is
$\kappa d-d/6=41d/96>2d/5$.
For II and III the nonexceptional bounds are
$\kappa d-d/12=49d/96>d/2$ and
$\kappa d-d/5=63d/160>3d/8$.
The exceptional coefficient in II and III can be raised to $h$ at
cost at most $(81d/160)\bar f_{L-1}$ or
$(59d/96)\bar f_{L-1}$. The actual bound $\bar f_{L-1}\le2$ pays
exactly the stated $E$, even if the index is outside the last four terms.
Retaining only those four actual terms then gives
\begin{equation}\tag{R20.8.1}\label{eq:R20:long}
e^\Gamma T^LF\le Y=1-h(4-d)+Q_{h,n}(B)u+\mu_\ell B^LR+E\ge0,
\quad Q_{h,n}=G_n-h(3+2z+z^2).
\end{equation}
The entire finite bound \eqref{eq:R18:shiftfour} gives
$\|Q_{h,n}\|\le1-3h$. Retaining $D_L$ in the actual mean yields
$m\le1-5h+3hd+E\le1-kd$ with
$k_I=37/20$, $k_{II}=13/10$, $k_{III}=97/192$.
The variance row bounds are $(17/15)^2d$, $(11/10)^2d$,
and $(7/6)^2d$, using only $L\ge1$, $\mu_\ell\le1$.

The pointwise form of \eqref{eq:R20:long} is
$Y=(1-3h)q+BV-h(1-d)+E$,
$V=[G_{n-1}-h(2+z)]u+\mu_\ell B^{L-1}R$.
Thus $\|V\|_2\le[1+\rho+h\sqrt{4+\theta}]\sqrt d$, with row
bounds $4/3$, $41/32$, $7/5$. They give respectively
\[
\begin{array}{c|l|c}
I&C_a+Y+m\le3-(21/5)d+(38/15)\sqrt d\le7/2&19/2\\
II&C_a+Y+m\le3-(109/40)d+(5/2)\sqrt d\le18/5&10\\
III&C_a+Y+m\le3-(77/48)d+(13/5)\sqrt d<4&13.
\end{array}
\]
The final column is the curvature coefficient. The first two full
increment maxima are $361/945<1/2$, $125/218<3/5$; the third is
at most $39/40<1$ using $\sqrt d<3/8$.
For $\phi(x)=1-3x+2x^2-x^3$,
$\phi(\theta)=1451/32768$,
$\operatorname{Re}(G_aG_n)\ge\phi$ and
$\operatorname{Re}[G_a(3+2z+z^2)]\le3$.
The shifted core has $\int Q_e\le d$ and
$\int Q_o=\beta_{n+1}-2hd\ge\tau-2hd>0$.
Its actual scalar correction is at most $d^2-\tau^2+2h\tau d$,
so
\[
\operatorname{Cov}(C_a,Q_{h,n}(B)u)
\le-\phi(\theta)d+[3(h/d)+1-\kappa^2]d^2+2(h/d)d^3.
\]
The normalized bound in I is $32629/163840<1/5$;
in II it is $7857/32768<1/4$; III is smaller than I.
Add the complete tail costs $3d/20$, $9d/80$, and $3d/16$,
respectively. The total covariance is at most $3d/8$ in I and II,
and $2d/5$ in III. The constant $E$ has no covariance.
The row mean exponents are $3677/160>39/2$, $633/32>39/2$,
and $14113/480>29$; the remainders are at most
\[
\begin{gathered}
(19/2)(289/225)+3/2=3083/225<53/3,\\
10(121/100)+3/2=68/5<53/3,\\
13(49/36)+8/5=3473/180<22.
\end{gathered}
\]
Since $\Gamma\ge16d/5$ in I,II and $\Gamma\ge32d/5$ in III,
\eqref{eq:R20:payment} proves the full long low range.

For $1/8\le d\le1/3$, the same actual envelope $|R|^2<d/3$ gives
$|R|\le(5/3)d$, $\Gamma\ge(16/5)jd$, and
$F\le e^{-(23/15)d}\bar f_\ell\le e^{-d}\bar f_\ell$.
Thus $T^LF\le e^{-d}\bar f_n$, $n=\ell+L\ge3$.
The three-step Module R2-C012, four-step Module R2-C014,
and $n\ge5$ Module R2-C017 give $\E[f_a\bar f_n^4]\le1$, so
\[
\E[f_a(T^LF)^4]\le e^{-4d}\le1-\frac{4d}{1+4d}
\le1-12d/7\le1-d/5.
\]
This joins at $d=1/8$. If the second shortest residual arm is three,
the shortest must be two and Module R2-C018 applies. Otherwise the
family above applies, proving the announced complete nonunit odd
terminal union. Directions use only this structural distinction.
At $d=0$ all actual messages are one by original leaf integration,
without implying $K=0$. Positivity restores all original equal pairs;
full reversal reverses the receiving arm, every terminal arm, each
connector edge, and all decorations simultaneously.
